%% Methods section for Nature paper

\subsection*{Unfolding and gap ratio computation}

Zeros $\gamma_n$ of $\zeta(\tfrac{1}{2} + i\gamma)$ are unfolded to unit mean
spacing via the smooth counting function~\eqref{eq:unfold}.
Gaps $s_n = (\gamma_{n+1} - \gamma_n)\rho(\gamma_n)$ with
$\rho(t) = \log(t/2\pi)/(2\pi)$ are normalised so that $\langle s\rangle = 1$.
Pathological gaps ($s < 0.02$ or $s > 6$) are excluded (fraction $< 0.1\%$).
The gap ratio $r_n = \min(s_n, s_{n+1})/\max(s_n, s_{n+1})$ is averaged over
each dataset.

\subsection*{Uncertainty estimation}

For the Platt data ($N = 500{,}000$ zeros each), the empirical standard error
is $\sigma_{\mathrm{emp}} = \mathrm{std}(r_n)/\sqrt{N_{\mathrm{pairs}}}$.
We impose a floor $\sigma_{\mathrm{floor}} = 0.00020$, corresponding to
$\sigma_{\mathrm{naive}}(N{=}500\mathrm{k})/\sqrt{10}$, to account for
residual correlations from the finite range of consecutive gaps used.
The final uncertainty is $\sigma = \max(\sigma_{\mathrm{emp}},\, \sigma_{\mathrm{floor}})$.

\subsection*{Model fitting}

All models are fit using weighted least squares with
$\chi^2 = \sum_i [(r_i - r_{\mathrm{model}}(\log T_i))/\sigma_i]^2$.
AIC is computed as $\chi^2 + 2k$ where $k$ is the number of parameters.
Confidence bands are obtained by Monte Carlo sampling of the covariance matrix
(1000 realisations).

\subsection*{Sensitivity analysis (CUE Monte Carlo)}

The sensitivities $\partial\langle r\rangle/\partial\,\mathrm{std}$ and
$\partial\langle r\rangle/\partial\,\mathrm{Corr}$ are measured from
Circular Unitary Ensemble (CUE) matrices of size $N = 2000$:
\begin{itemize}
\item Method~A (std): rescale $s' = 1 + \lambda(s-1)$ to change std without
changing correlation.
\item Method~B (Corr): shuffle a fraction of $s_2$ values to change correlation
without changing std.
\end{itemize}

\subsection*{Fredholm determinant and $p_2^{\mathrm{GUE}}$}

The gap probability $E(s) = \det(I - K_s)$ is computed via Gauss--Legendre
quadrature with $n_{\mathrm{quad}} = 32$ nodes~\cite{bornemann2010}.
The joint density of two consecutive gaps $p_2(s_1, s_2)$ uses the
three-point Schur complement formula:
\begin{equation}
p_2(s_1, s_2) = \det_3(0, s_1, S) \times E^{(0, s_1, S)}(S)
\end{equation}
where $E^{(0, s_1, S)}$ is the Fredholm determinant of the conditional kernel
obtained by Schur complement of the $3\times 3$ Gram matrix.

\subsection*{Conrey--Snaith triple correlation}

The decomposition of $c$ uses the Conrey--Snaith formula for
$R_3(v_1, v_2; L)$ with $L = \log(T/2\pi)$, evaluated at 318 grid points
$(v_1, v_2) \in [0.3, 2.5]^2$ with step $0.1$ and
$L \in \{30, 40, 50, 60, 80, 1000, 3000, 10{,}000, 30{,}000\}$.
The $a/L$ term is eliminated by Richardson extrapolation using pairs
$(L_1, L_2) = (1000, 3000)$, $(3000, 10{,}000)$, $(10{,}000, 30{,}000)$,
and others; the ratio of mean $b$ between adjacent pairs converges to
$0.979$ with inter-pair correlation $r = 0.9999$.
The $R_3$ channel contribution is $c_{R_3} = \sum (r - R_{\mathrm{GUE}}) \cdot b
\cdot E^{\mathrm{cond}}\, (\Delta v)^2 / Z_{p_2}$, where $E^{\mathrm{cond}}$ is
the conditional gap probability from the Schur complement of the sine kernel
and $r(v_1, v_2) = \min(s_1, s_2)/\max(s_1, s_2)$ with $s_1 = \min(v_1, v_2)$,
$s_2 = |v_1 - v_2|$ (positions convention). This yields $c_{R_3} = -1.25$;
the gap probability channel $c_E = +2.50$ follows by difference.
